% Sources: README "Viscous torque and disc temperature", "Luminosity budget"
% Code: ULXspectra, branch fix-results-integrity, commit 5f00569

\section{Disc structure and luminosity budget}
\label{sec:disc}

\subsection{Viscous torque}
\label{sec:disc:torque}

Let $h = \sqrt{GMR}$ be the specific angular momentum of the Keplerian disc and $\Fvis > 0$ the viscous torque.
If the matter lost to the wind at $R$ carries away the specific angular momentum $k_{\rm w}h(R)$, the conservation of angular momentum in a stationary disc reads
\begin{equation}
\frac{{\rm d}(\Mdot h)}{{\rm d}R} - \frac{{\rm d}\Fvis}{{\rm d}R} - k_{\rm w}\,h\,\frac{{\rm d}\Mdot}{{\rm d}R} = 0 .
\label{eq:angmom}
\end{equation}
% source: old manuscript (SEA_spectrum.zip, main.tex, Sect. "Effective temperature distribution"),
%         there written for F = -F_vis; code: disc_Fvis_point() (csrcs/disc_struct.c) assumes k_w = 1
We take $k_{\rm w} = 1$, that is, the wind removes the Keplerian angular momentum of the radius where it is launched.
Then the viscous torque obeys
\begin{equation}
\frac{{\rm d}\Fvis}{{\rm d}h} = \Mdot(R), \quad \Fvis(R) = F_{\rm in} + \int_{h(\rin)}^{h(R)} \!\Mdot\,{\rm d}h .
\label{eq:Fvis}
\end{equation}
% code: disc_Fvis_point() (csrcs/disc_struct.c); README "Viscous torque and disc temperature"; PLAN C10
For a BH the torque vanishes at the inner radius, $F_{\rm in} = 0$.
For a NS, the magnetosphere exerts a torque at $\rmag$,
\begin{equation}
F_{\rm in} = \beta_{\rm t}\,\Mdot(\rmag)\,h(\rmag),
\label{eq:Fin_NS}
\end{equation}
% code: disc_Fvis_point() (csrcs/disc_struct.c); README "Viscous torque and disc temperature"
with $\beta_{\rm t} = 1$.
% AUTHOR: physical justification of beta_t = 1 (PLAN C10)?

All the profiles of Sect.~\ref{sec:model:mdot} are linear in $R$ between $\rin$ and $\Rsph$, so the integral is analytic.
With $s = \sqrt{\min(x, \xsp)}$ and $h_{\rm in} = h(\rin)$,
\begin{multline}
\Fvis = F_{\rm in} + \frac{h_{\rm in}(s - 1)}{3(s + 1)}\bigl[\Mdot(\rin)(2s + 1) \\
+ \Mdot(s^2\rin)(s + 2)\bigr]
+ \Mdot_0\,\bigl[h(R) - h(\Rsph)\bigr]\,\Theta(R - \Rsph),
\label{eq:Fvis_lin}
\end{multline}
% code: disc_Fvis_point() (csrcs/disc_struct.c), closed form in its header comment
where $\Theta$ is the Heaviside step function.
$\Fvis$ and the disc temperature are continuous at $\Rsph$.

Since the wind takes away angular momentum, less of it has to be transported by the viscous stress, and it is the stress that heats the disc.
Outside $\Rsph$ the torque is lower than in a disc without outflow and with the same $\Mdot_0$ by $\int (\Mdot_0 - \Mdot)\,{\rm d}h$, taken from $\rin$ to $\Rsph$, so the outer disc is colder.
% source: old manuscript (SEA_spectrum.zip, main.tex), paragraph "Worth noting that the presence of the wind makes the disc colder"
% derived (Claude): the difference follows from eq. (Fvis) with F_in equal in both discs.
% AUTHOR: the old manuscript marks the paragraph before this one "somehow wrong: see the
%         figure" (Fvis_comp.png: SS73 torque vs. an advective disc with outflow). Is the
%         statement above unaffected?
The relative difference fades at $R \gg \Rsph$, where $\Fvis \approx \Mdot_0 h$ in both discs.
% derived (Claude)
% AUTHOR (minor): for P07 with r_in,f < 6 the profile can be clipped at 0 near r_in;
%         the code ignores that non-linear piece (comment in disc_Fvis_point()).
%         Mention or not?

\subsection{Effective temperature}
\label{sec:disc:teff}

The viscous heating per unit area of one side of the disc is
\begin{equation}
\Qvis = \frac{3}{8\pi}\,\frac{\Fvis\,\Omega_{\rm K}}{R^2} .
\label{eq:Qvis}
\end{equation}
% code: disc_Qvis_point() (csrcs/disc_struct.c); README "Viscous torque and disc temperature"
Outside the trapping radius the disc radiates locally, $\sigma\Teff^4 = \Qvis$.
Inside it, advection makes the radiated flux fall as $R^{-2}$, and we join the two continuously at $\Rtrap$:
\begin{equation}
\sigma\Teff^4 =
\begin{cases}
\Qvis(R), & R > \Rtrap, \\
\Qvis(\Rtrap)\,(\Rtrap/R)^2, & \rin \le R \le \Rtrap .
\end{cases}
\label{eq:Teff}
\end{equation}
% code: disc_Qvis_and_Teff() (csrcs/disc_struct.c); README "Viscous torque and disc temperature"
$\Qvis(\Rtrap)$ is evaluated as if the disc continued beyond $\rout$.
% AUTHOR: physical argument for sigma T^4 ~ R^-2 inside R_trap (L_per ln R = const),
%         with a reference?

\subsection{Luminosity budget}
\label{sec:disc:budget}

The disc luminosity is the integral of the profile~\eqref{eq:Teff} over both sides of the disc from $\rin$ to $\rout$:
\begin{multline}
L_{\rm disc} = 4\pi\int_{\rin}^{\rout} \sigma\Teff^4\,R\,{\rm d}R
= L_{\rm vis}(\Rtrap, \rout) \\
+ 4\pi\Rtrap^2\,\Qvis(\Rtrap)\,\ln\frac{\min(\Rtrap, \rout)}{\rin},
\label{eq:Ldisc}
\end{multline}
% code: accretion_luminosity_split(), disc_L_out_above(), disc_advective_L_per_lnr() (csrcs/wind_emission.c, csrcs/disc_struct.c); README "Luminosity budget"
where $L_{\rm vis}(\Rtrap, \rout) = 4\pi\int \Qvis R\,{\rm d}R$ is the viscous release between $\Rtrap$ and $\rout$.
For $\Rtrap \le \rin$ this is the whole viscous release, which also covers sub-Eddington rates.
When $\Rtrap$ is close to the radius outside which the disc releases $\LEdd$, Eq.~\eqref{eq:Ldisc} gives $L_{\rm disc} \approx \LEdd[1 + \ln(\Rtrap/\rin)]$.
% code: README "Luminosity budget", mode 1+log

For a BH, nothing inside $\rin$ radiates.
For a NS, the matter reaching $\rmag$ falls onto the star along the field lines and releases
\begin{equation}
L_{\rm central} = \eta_{\rm magn}\,G M\,\Mdot(\rmag)\left(\frac{1}{\RNS} - \frac{1}{\rmag}\right),
\label{eq:Lcentral}
\end{equation}
% code: accretion_luminosity_split(), mode "NS" (csrcs/wind_emission.c); README "Luminosity budget"
with the efficiency $\eta_{\rm magn} \le 1$, which covers losses such as outflows from the magnetosphere.
The Keplerian kinetic energy of the matter at $\rmag$, about $GM\Mdot/(2\rmag)$, is not counted: for a slow rotator it would also be released inside, while near spin equilibrium it goes partly into the spin of the star.
% code: comment at accretion_luminosity_split() (csrcs/wind_emission.c)
% AUTHOR: check that this is how you want the omission justified.
The total accretion luminosity is $L_{\rm acc} = L_{\rm disc} + L_{\rm central}$.
% code: diagnostic Lbol_accr_teor
